\chapter{移位寄存器构建环形计数器}

\section{从移位寄存器出发}

回顾移位寄存器：数据从一端进入，经过N拍从另一端出来。

现在思考一个问题：如果把串行输出\textbf{反馈}到串行输入呢？

\begin{center}
\begin{tikzpicture}
  \node[draw, minimum width=1.2cm, minimum height=1.5cm] (dff0) at (0,0) {DFF0};
  \node[draw, minimum width=1.2cm, minimum height=1.5cm] (dff1) at (2,0) {DFF1};
  \node[draw, minimum width=1.2cm, minimum height=1.5cm] (dff2) at (4,0) {DFF2};
  \node[draw, minimum width=1.2cm, minimum height=1.5cm] (dff3) at (6,0) {DFF3};

  \draw[->] (dff0) -- (dff1);
  \draw[->] (dff1) -- (dff2);
  \draw[->] (dff2) -- (dff3);
  % Feedback
  \draw[->] (dff3.east) -- ++(0.5,0) |- ++(0,-1.5) -- ++(-6.5,0) |- (dff0.west);
  \node at (3,-2.2) {反馈环：Q3 → D0};
\end{tikzpicture}
\end{center}

\begin{keypoint}[环形计数器的形成]
\textbf{环形计数器 = 移位寄存器 + Q输出反馈到串行输入}

数据不再从外部输入，而是在D触发器链中\textbf{循环}流动。

$$\text{数据在闭合的环中无限循环}$$
\end{keypoint}

\section{逐拍分析：环形计数器如何工作}

假设4位环形计数器，初始状态通过复位设为 1000（即Q3=1, Q2=0, Q1=0, Q0=0）。

注意：这里假设Q3是最高位（最左），Q0是最低位（最右）。

反馈连接：Q0 → D3（最低位反馈回最高位的输入）。

\subsection{时钟第0拍（初始/复位后）}

寄存器内容：[Q3 Q2 Q1 Q0] = [1 0 0 0]

\subsection{时钟第1拍上升沿}

\begin{beatanalysis}[环形计数器——时钟第1拍]
\begin{itemize}
  \item D3 = Q0 = 0（反馈值）
  \item D2 = Q3 = 1
  \item D1 = Q2 = 0
  \item D0 = Q1 = 0
\end{itemize}

\textbf{上升沿后}：
\begin{itemize}
  \item Q3 ← D3 = 0（从Q0反馈回来的0）
  \item Q2 ← D2 = 1（Q3原来的1右移一位）
  \item Q1 ← D1 = 0
  \item Q0 ← D0 = 0
\end{itemize}

\textbf{寄存器内容：[0 1 0 0]}（1向右移动了一位！）
\end{beatanalysis}

\subsection{时钟第2拍上升沿}

\begin{beatanalysis}[环形计数器——时钟第2拍]
\begin{itemize}
  \item D3 = Q0 = 0
  \item D2 = Q3 = 0
  \item D1 = Q2 = 1
  \item D0 = Q1 = 0
\end{itemize}

\textbf{上升沿后：[0 0 1 0]}（1继续右移）
\end{beatanalysis}

\subsection{时钟第3拍上升沿}

\begin{beatanalysis}[环形计数器——时钟第3拍]
\begin{itemize}
  \item D3 = Q0 = 0
  \item D2 = Q3 = 0
  \item D1 = Q2 = 0
  \item D0 = Q1 = 1
\end{itemize}

\textbf{上升沿后：[0 0 0 1]}（1移到了最右边）
\end{beatanalysis}

\subsection{时钟第4拍上升沿}

\begin{beatanalysis}[环形计数器——时钟第4拍：关键时刻！]
\begin{itemize}
  \item D3 = Q0 = \textbf{1}（反馈环起作用！Q0的1反馈回D3）
  \item D2 = Q3 = 0
  \item D1 = Q2 = 0
  \item D0 = Q1 = 0
\end{itemize}

\textbf{上升沿后：[1 0 0 0]}（1回到最左边——完成一圈循环！）
\end{beatanalysis}

\subsection{完整的循环过程}

\begin{center}
\begin{tabular}{|c|c|c|c|c|}
\hline
\textbf{时钟拍} & \textbf{Q3} & \textbf{Q2} & \textbf{Q1} & \textbf{Q0} \\
\hline
0（初始）& \textbf{1} & 0 & 0 & 0 \\
1 & 0 & \textbf{1} & 0 & 0 \\
2 & 0 & 0 & \textbf{1} & 0 \\
3 & 0 & 0 & 0 & \textbf{1} \\
4 & \textbf{1} & 0 & 0 & 0 \\
5 & 0 & \textbf{1} & 0 & 0 \\
... & ... & ... & ... & ... \\
\hline
\end{tabular}
\end{center}

\textbf{这个"1在循环移动"的现象就是环形计数器的本质。}

\section{环形计数器 vs 普通计数器}

\begin{center}
\begin{tabular}{|l|c|c|}
\hline
\textbf{特性} & \textbf{二进制计数器} & \textbf{环形计数器} \\
\hline
硬件结构 & 加法器 + D触发器 & N个D触发器 + 反馈环 \\
状态编码 & 二进制（000,001,...）& One-hot（1000,0100,...）\\
N个状态需要 & $\lceil\log_2 N\rceil$个DFF & N个DFF \\
状态解码 & 需要组合逻辑 & 直接看哪个DFF的Q=1 \\
速度 & 受加法器进位链限制 & 最快（只有移位） \\
最大状态数 & $2^k$ & k（k个DFF产生k个状态） \\
\hline
\end{tabular}
\end{center}

环形计数器用更多触发器换取更快的速度和更简单的解码。

\section{为什么反馈后形成循环状态}

反馈把串行输出接回串行输入，形成一个\textbf{闭合回路}。

\begin{itemize}
  \item 没有反馈：数据从一端进入，从另一端消失（一去不回）
  \item 有反馈：数据从一端出来后，重新从另一端进入（无限循环）
\end{itemize}

这就像一条首尾相接的传送带——东西在上面无限循环移动。

\section{约翰逊计数器（扭环形计数器）}

约翰逊计数器是环形计数器的变体：将\textbf{反相的}Q输出反馈回输入。

\begin{itemize}
  \item 环形计数器：Q0 → D3（同相反馈）
  \item 约翰逊计数器（扭环形）：$\overline{Q_0}$ → D3（反相反馈）
\end{itemize}

\textbf{效果：k个触发器可以产生2k个状态}（比普通环形计数器多一倍！）

4位约翰逊计数器的状态序列：

$$0000 \to 1000 \to 1100 \to 1110 \to 1111 \to 0111 \to 0011 \to 0001 \to 0000$$

共8个状态（4个触发器产生8个状态）。

\section{VHDL实现}

\begin{lstlisting}[caption={4位环形计数器}]
library ieee;
use ieee.std_logic_1164.all;

entity ring_counter is
  port (
    clk   : in  std_logic;
    rst_n : in  std_logic;
    q     : out std_logic_vector(3 downto 0)
  );
end ring_counter;

architecture behavioral of ring_counter is
  signal reg : std_logic_vector(3 downto 0);
begin
  process(clk, rst_n)
  begin
    if rst_n = '0' then
      reg <= "1000";  -- 初始状态：只有最高位为1
    elsif rising_edge(clk) then
      reg <= reg(0) & reg(3 downto 1);  -- 右移 + 最低位反馈到最高位
    end if;
  end process;

  q <= reg;
end behavioral;
\end{lstlisting}

\textbf{关键语句}：\texttt{reg <= reg(0) \& reg(3 downto 1)}

这是右移 + 反馈的合并：
\begin{itemize}
  \item reg(0)（最低位）$\to$ reg(3)（反馈到最高位输入）
  \item reg(3) $\to$ reg(2)（右移）
  \item reg(2) $\to$ reg(1)（右移）
  \item reg(1) $\to$ reg(0)（右移）
\end{itemize}

\subsection{约翰逊计数器VHDL}

\begin{lstlisting}[caption={4位约翰逊计数器（扭环形）}]
architecture behavioral of johnson_counter is
  signal reg : std_logic_vector(3 downto 0);
begin
  process(clk, rst_n)
  begin
    if rst_n = '0' then
      reg <= (others => '0');  -- 初始全0
    elsif rising_edge(clk) then
      reg <= (not reg(0)) & reg(3 downto 1);  -- 反相的最低位移入最高位
    end if;
  end process;
end behavioral;
\end{lstlisting}

\section{考试常见变式}

\begin{enumerate}
  \item \textbf{变式1：设计N位环形计数器}

  N位环形计数器产生N个状态，初始值100...0（one-hot编码）。

  \item \textbf{变式2：设计2N位约翰逊计数器}

  N个触发器产生2N个状态。反馈是反相的最低位移入最高位。

  \item \textbf{变式3：自启动的环形计数器}

  如果因干扰进入非法状态（如0101，两个1），需要自动回到合法状态。这需要额外的修正逻辑。

  \item \textbf{变式4：比较环形计数器与二进制计数器}

  常考：N状态各需多少触发器？各有哪些优缺点？
\end{enumerate}

\section{Testbench}

\begin{lstlisting}[caption={环形计数器Testbench}]
library ieee;
use ieee.std_logic_1164.all;

entity ring_counter_tb is
end ring_counter_tb;

architecture test of ring_counter_tb is
  signal clk, rst_n : std_logic;
  signal q : std_logic_vector(3 downto 0);
  constant clk_period : time := 20 ns;

  component ring_counter is
    port (clk, rst_n : in std_logic; q : out std_logic_vector(3 downto 0));
  end component;
begin
  uut: ring_counter port map (clk, rst_n, q);

  clk_process: process
  begin
    clk <= '0'; wait for clk_period/2;
    clk <= '1'; wait for clk_period/2;
  end process;

  stimulus: process
  begin
    rst_n <= '0'; wait for 30 ns;
    rst_n <= '1';
    -- 观察两个完整循环（8拍）
    wait for clk_period * 9;
    wait;
  end process;
end test;
\end{lstlisting}

\begin{examtip}[环形计数器秒杀法]
\textbf{初始}：reg = "1000"（只有一个1）

\textbf{每个时钟沿}：1向右移动一位（因为反馈环）

\textbf{核心VHDL}：\texttt{reg <= reg(0) \& reg(N-1 downto 1)}

\textbf{约翰逊计数器}：用反相：\texttt{reg <= (not reg(0)) \& reg(N-1 downto 1)}
\end{examtip}

\section{变式详解}
\chapter{环形计数器——全部变式详解}

\section{变式1：N位环形计数器}

\begin{detailbox}[完整教学]
\textbf{【通用模板】}
\begin{lstlisting}
process(clk, rst_n)
begin
  if rst_n='0' then reg <= (N-1=>'1', others=>'0');  -- One-hot初值
  elsif rising_edge(clk) then
    reg <= reg(0) & reg(N-1 downto 1);  -- 右移+反馈
  end if;
end process;
\end{lstlisting}
\textbf{【结论】}N个DFF→N个状态。所有状态中恰好一个1（One-hot编码）。
\end{detailbox}


\section{变式2：约翰逊计数器（2N个状态）}

\begin{detailbox}[完整教学]
\textbf{【反馈】}\texttt{reg <= (not reg(0)) \& reg(N-1 downto 1);}
\textbf{【4位约翰逊状态序列】}$0000\to1000\to1100\to1110\to1111\to0111\to0011\to0001\to0000$
\textbf{【结论】}N个DFF→2N个状态（比环形计数器多一倍）。
\end{detailbox}


\section{变式3：自启动环形计数器}

\begin{detailbox}[完整教学]
\textbf{【问题】}干扰可能使环形计数器进入非法状态。
\textbf{【方案】}检测非法状态（如全0或多个1），强制恢复到合法状态(1000)。
\textbf{【VHDL】}在移位前检查：\texttt{if count\_ones(reg) /= 1 then reg <= "1000";}
\end{detailbox}


\section{变式4：用环形计数器实现时序控制器}

\begin{detailbox}[完整教学]
\textbf{【应用】}多周期操作器的步骤控制。
\textbf{【5步控制器】}5位环形计数器→每拍恰好一个步骤信号为1→控制对应步骤。
\textbf{【优点】}无需额外的译码逻辑——直接读Q值。
\end{detailbox}
